\documentclass[11pt]{article}
\usepackage[margin=1in]{geometry}
\usepackage{amsmath,amssymb,amsthm,mathtools}
\usepackage{hyperref}

\newtheorem{theorem}{Theorem}
\newtheorem{lemma}[theorem]{Lemma}
\newtheorem{proposition}[theorem]{Proposition}
\newtheorem{corollary}[theorem]{Corollary}
\theoremstyle{definition}
\newtheorem{definition}[theorem]{Definition}
\theoremstyle{remark}
\newtheorem{remark}[theorem]{Remark}

\DeclareMathOperator{\ImOp}{Im}
\DeclareMathOperator{\ReOp}{Re}
\DeclareMathOperator{\supp}{supp}
\let\Im\ImOp
\let\Re\ReOp

\title{The Complex Orthogonal Measure\\
and the Spectral Representation\\
of the Hardy \texorpdfstring{$Z$}{Z}-Function}
\author{}
\date{}

\begin{document}
\maketitle

\begin{abstract}
Let $\theta$ be the Riemann--Siegel theta function, $Z(t)=e^{i\theta(t)}\zeta(\tfrac12+it)$ the Hardy $Z$-function, and $H(x):=Z(\theta^{-1}(x))/\sqrt{\theta'(\theta^{-1}(x))}$ the warped Hardy signal. Given as input the companion results
\begin{itemize}
\item Theorem~A of \emph{Three Missing Theorems}, which establishes
\[
 C(\eta)\;:=\;\lim_{X\to\infty}\frac{1}{X-X_0}\int_{X_0}^{X-|\eta|} H(x)\overline{H(x+|\eta|)}\,dx\;=\;2\cos|\eta|,\qquad \eta\in\mathbb{R},
\]
with the quantitative rate $|C_X(\eta)-C_{\mathrm{diag}}(\eta)|\le 8\pi(1+|\eta|)\log^3(X/X_0)/(X-X_0)+0.349/\sqrt{T_0}$;
\item the Bochner factorization $C(\eta)=\int_{\mathbb{R}} e^{i\eta\xi}\,d\nu(\xi)$ with $\nu=\delta_{-1}+\delta_{+1}$ from \emph{Existence of the Spectral Measure in the Limit},
\end{itemize}
we define the complex orthogonal measure $\Phi$ of $H$ by an \emph{explicit Cram\'er-type pre-limit formula}, prove convergence of the pre-limit in the time-average Hilbert space generated by the shifts of $H$, and verify by Fubini the spectral representation
\[
  H(y)\;=\;\int_{\mathbb{R}} e^{i\xi y}\,d\Phi(\xi),\qquad y\in\mathbb{R},
\]
as an identity in that Hilbert space. The measure $\Phi$ is Hilbert-space-valued, orthogonally-scattered on the Borel $\sigma$-algebra of $\mathbb{R}$, supported on $\{-1,+1\}$, with intensity $\langle\Phi(A),\Phi(B)\rangle=\nu(A\cap B)$. The corresponding representation of the Hardy $Z$-function is
\[
  Z(t)\;=\;\sqrt{\theta'(t)}\int_{\mathbb{R}} e^{i\xi\theta(t)}\,d\Phi(\xi)
\]
in the same Hilbert space.
\end{abstract}

\section{Inputs and notation}\label{sec:inputs}

Fix $T_0\ge 200$ and set $X_0:=\theta(T_0)$. For $X>X_0$ set $T:=\theta^{-1}(X)$. The Riemann--Siegel theta function $\theta:[T_0,\infty)\to[X_0,\infty)$ is a diffeomorphism with $\theta'(t)>0$, and $H\in C_{\mathrm{loc}}^\infty([X_0,\infty);\mathbb{C})$ is locally bounded.

\paragraph{Time-average inner product.} For $f,g\colon[X_0,\infty)\to\mathbb{C}$ bounded and measurable, write, whenever the limit exists,
\[
  \langle f,g\rangle_\sharp\;:=\;\lim_{X\to\infty}\frac{1}{X-X_0}\int_{X_0}^X f(x)\overline{g(x)}\,dx.
\]
Theorem~A states exactly that the limit $\langle H_s,H_t\rangle_\sharp$ exists for every $s,t\in\mathbb{R}$, where
\[
  H_s(x)\;:=\;H(x+s)\cdot \mathbf{1}_{\{x+s\ge X_0\}},
\]
and equals $C(s-t)=2\cos(s-t)=e^{-i(s-t)}+e^{+i(s-t)}$. (The indicator does not affect the time average, which is insensitive to a compact perturbation of the integrand.)

\paragraph{Pre-Hilbert space.} Let $\mathcal{E}$ denote the linear span of $\{H_s:s\in\mathbb{R}\}$ over $\mathbb{C}$. Define the sesquilinear form on $\mathcal{E}$ by
\[
  \Big\langle\sum_k\alpha_k H_{s_k},\sum_j\beta_j H_{t_j}\Big\rangle_\sharp\;:=\;\sum_{k,j}\alpha_k\overline{\beta_j}\,C(s_k-t_j).
\]
This is well-defined because each individual time average $\langle H_{s_k},H_{t_j}\rangle_\sharp$ exists by Theorem~A and is bilinear in the coefficients. The Gram matrix $[C(s_k-t_j)]_{k,j}$ is Hermitian and positive semidefinite: by the Bochner factorization $\nu=\delta_{-1}+\delta_{+1}$,
\[
  \sum_{k,j}\alpha_k\overline{\alpha_j}\,C(s_k-s_j)\;=\;\int_{\mathbb{R}}\Big|\sum_k\alpha_k e^{i\xi s_k}\Big|^2 d\nu(\xi)\;=\;\Big|\sum_k\alpha_k e^{-is_k}\Big|^2+\Big|\sum_k\alpha_k e^{+is_k}\Big|^2\;\ge\;0.
\]
Quotient $\mathcal{E}$ by its null space $\mathcal{N}=\{f\in\mathcal{E}:\langle f,f\rangle_\sharp=0\}$ and complete in the resulting norm. Call the resulting Hilbert space $\mathcal{H}$. The natural injection $\iota:\mathcal{E}/\mathcal{N}\hookrightarrow\mathcal{H}$ has dense image.

Write $\widetilde{H_s}\in\mathcal{H}$ for the image of $H_s\in\mathcal{E}$ under the composition $\mathcal{E}\to\mathcal{E}/\mathcal{N}\hookrightarrow\mathcal{H}$. Then
\[
  \langle\widetilde{H_s},\widetilde{H_t}\rangle_{\mathcal{H}}\;=\;C(s-t)\;=\;e^{-i(s-t)}+e^{+i(s-t)}.
\]

\paragraph{Extension of the inner product to Bochner integrals of $\mathcal{H}$-valued functions.} For a continuous compactly supported function $K:[X_0,\infty)\to\mathbb{C}$, the $\mathcal{H}$-valued integral
\[
  \int_{X_0}^\infty K(x)\,\widetilde{H_x}\,dx\;\in\;\mathcal{H}
\]
is defined as the Bochner integral of the continuous $\mathcal{H}$-valued map $x\mapsto K(x)\widetilde{H_x}$. (Continuity in $x$ in the $\mathcal{H}$-norm is implied by $\|\widetilde{H_{x+h}}-\widetilde{H_x}\|_{\mathcal{H}}^2=C(0)-2\Re C(h)+C(0)=4-4\cos h\to 0$ as $h\to 0$.) Its squared norm satisfies the usual Bochner identity
\[
  \Big\|\int K(x)\widetilde{H_x}\,dx\Big\|_{\mathcal{H}}^2\;=\;\iint K(x)\overline{K(y)}\,C(x-y)\,dx\,dy.
\]

\section{The explicit pre-limit Cram\'er formula}\label{sec:prelimit}

\begin{definition}[Pre-limit band coefficients]\label{def:prelimit-band}
For $X>X_0$ set
\[
  \Phi_{+,X}\;:=\;\frac{1}{X-X_0}\int_{X_0}^X e^{-ix}\,\widetilde{H_x}\,dx\;\in\;\mathcal{H},\qquad
  \Phi_{-,X}\;:=\;\frac{1}{X-X_0}\int_{X_0}^X e^{+ix}\,\widetilde{H_x}\,dx\;\in\;\mathcal{H}.
\]
\end{definition}

The two integrals are Bochner integrals of the continuous $\mathcal{H}$-valued integrands $x\mapsto e^{\mp ix}\widetilde{H_x}$, as in the paragraph closing Section~\ref{sec:inputs}. They are explicitly given by a finite Cesaro average in $\mathcal{H}$ and are defined for every $X>X_0$, without passing to any limit.

\begin{definition}[Pre-limit orthogonal measure]\label{def:prelimit-measure}
For every Borel set $A\subset\mathbb{R}$ and every $X>X_0$ define
\[
  \Phi_X(A)\;:=\;\mathbf{1}_A(+1)\,\Phi_{+,X}\;+\;\mathbf{1}_A(-1)\,\Phi_{-,X}\;\in\;\mathcal{H}.
\]
\end{definition}

$\Phi_X(\cdot)$ is a finitely supported $\mathcal{H}$-valued set function carried on $\{-1,+1\}$. Countable additivity in $A$ is immediate from the two-atom form. The dependence on $X$ is entirely in the two coefficients $\Phi_{\pm,X}$.

\begin{remark}\label{rem:cramer-form}
Definition~\ref{def:prelimit-band} is the Cram\'er pre-limit formula specialised to the spectrum $\{\pm 1\}$ established by Theorem~A. For a general weakly harmonizable integrand with spectral measure $\nu$ on $\mathbb{R}$ and shift family $\widetilde{H_x}$, the Cram\'er pre-limit for a Borel set $A$ takes the form
\[
  \Phi_X(A)\;=\;\frac{1}{X-X_0}\int_{X_0}^X\widehat{\mathbf{1}_A}(x)\,\widetilde{H_x}\,dx,
\]
understood in a suitable regularised sense when $\widehat{\mathbf{1}_A}$ is not a tempered function. In the present case $\supp\nu=\{\pm 1\}$, so only the values of $\widehat{\mathbf{1}_A}$ at the two atoms contribute: $\widehat{\mathbf{1}_A}(x)$ is replaced by $\mathbf{1}_A(+1)e^{-ix}+\mathbf{1}_A(-1)e^{+ix}$, which is the formula in Definition~\ref{def:prelimit-band}.
\end{remark}

\section{Cauchy convergence of the pre-limit in \texorpdfstring{$\mathcal{H}$}{H}}\label{sec:convergence}

\begin{lemma}[Inner products of the pre-limit band coefficients]\label{lem:pre-inner-products}
For every $X,X'>X_0$,
\begin{align*}
  \langle \Phi_{+,X},\widetilde{H_y}\rangle_{\mathcal{H}} &\;=\;\frac{1}{X-X_0}\int_{X_0}^X e^{-ix}\,C(x-y)\,dx\;=\;e^{-iy}+R_+(X;y),\\[2pt]
  \langle \Phi_{-,X},\widetilde{H_y}\rangle_{\mathcal{H}} &\;=\;\frac{1}{X-X_0}\int_{X_0}^X e^{+ix}\,C(x-y)\,dx\;=\;e^{+iy}+R_-(X;y),
\end{align*}
where
\[
  R_+(X;y)\;=\;\frac{e^{iy}}{X-X_0}\int_{X_0}^X e^{-2ix}dx\;=\;\frac{e^{iy}}{-2i(X-X_0)}\left(e^{-2iX}-e^{-2iX_0}\right),
\]
and $R_-$ is its complex conjugate expression with the roles of the two exponents interchanged. In particular $|R_\pm(X;y)|\le 1/(X-X_0)$.
\end{lemma}

\begin{proof}
By the definition of the Bochner integral, $\langle\int K(x)\widetilde{H_x}dx,\widetilde{H_y}\rangle_\mathcal{H}=\int K(x)\langle\widetilde{H_x},\widetilde{H_y}\rangle_\mathcal{H}dx=\int K(x)C(x-y)dx$. Take $K(x)=e^{-ix}/(X-X_0)$ on $[X_0,X]$ and zero elsewhere. Since $C(x-y)=e^{-i(x-y)}+e^{+i(x-y)}$,
\[
  \int_{X_0}^X e^{-ix}C(x-y)dx\;=\;e^{iy}\int_{X_0}^X e^{-2ix}dx\;+\;e^{-iy}\int_{X_0}^X 1\,dx.
\]
Dividing by $X-X_0$ gives $e^{-iy}+R_+(X;y)$ with $R_+$ as stated. The oscillatory integral $\int_{X_0}^X e^{-2ix}dx$ is bounded in modulus by $1$, giving $|R_+(X;y)|\le 1/(X-X_0)$. The computation for $\Phi_{-,X}$ is identical with $-i$ replaced by $+i$ in the kernel.
\end{proof}

\begin{lemma}[Inner products of the pre-limit band coefficients among themselves]\label{lem:pre-self-inner}
For every $X,X'>X_0$,
\begin{align*}
  \langle\Phi_{+,X},\Phi_{+,X'}\rangle_{\mathcal{H}} &\;=\;1+E_{++}(X,X'),\\
  \langle\Phi_{-,X},\Phi_{-,X'}\rangle_{\mathcal{H}} &\;=\;1+E_{--}(X,X'),\\
  \langle\Phi_{+,X},\Phi_{-,X'}\rangle_{\mathcal{H}} &\;=\;E_{+-}(X,X'),
\end{align*}
with $|E_{\cdot\cdot}(X,X')|\le 2/\min(X-X_0,X'-X_0)$.
\end{lemma}

\begin{proof}
For the first inner product, using the Bochner identity,
\begin{align*}
  \langle\Phi_{+,X},\Phi_{+,X'}\rangle_{\mathcal{H}} &\;=\;\frac{1}{(X-X_0)(X'-X_0)}\int_{X_0}^X\int_{X_0}^{X'} e^{-ix}\,e^{+iy}\,C(x-y)\,dx\,dy\\
  &\;=\;\frac{1}{(X-X_0)(X'-X_0)}\int_{X_0}^X\int_{X_0}^{X'} e^{-ix+iy}\bigl[e^{-i(x-y)}+e^{+i(x-y)}\bigr]dx\,dy\\
  &\;=\;\frac{1}{(X-X_0)(X'-X_0)}\int_{X_0}^X\int_{X_0}^{X'}\bigl[e^{-2ix+2iy}+1\bigr]dx\,dy\\
  &\;=\;1\;+\;\frac{1}{(X-X_0)(X'-X_0)}\int_{X_0}^X e^{-2ix}dx\cdot\int_{X_0}^{X'}e^{+2iy}dy.
\end{align*}
Each one-variable oscillatory integral has modulus bounded by $1$, so the product of the two quotients has modulus bounded by $1/((X-X_0)(X'-X_0))$. Absorbing this into the stated bound $2/\min(X-X_0,X'-X_0)$ (valid whenever $X-X_0\ge 1$ or $X'-X_0\ge 1$) gives the claim. The other two computations are identical with signs interchanged; for the cross term,
\[
  \langle\Phi_{+,X},\Phi_{-,X'}\rangle_{\mathcal{H}}\;=\;\frac{1}{(X-X_0)(X'-X_0)}\int_{X_0}^X\int_{X_0}^{X'}\bigl[e^{-2ix}+e^{+2iy}\bigr]dx\,dy,
\]
which is bounded in modulus by the same quantity.
\end{proof}

\begin{proposition}[Cauchy convergence]\label{prop:cauchy}
The families $(\Phi_{+,X})_{X>X_0}$ and $(\Phi_{-,X})_{X>X_0}$ are Cauchy in $\mathcal{H}$ as $X\to\infty$. Their limits, denoted $\Phi_+:=\lim\Phi_{+,X}$ and $\Phi_-:=\lim\Phi_{-,X}$, exist in $\mathcal{H}$ and satisfy
\[
  \langle\Phi_+,\Phi_+\rangle_\mathcal{H}=\langle\Phi_-,\Phi_-\rangle_\mathcal{H}=1,\qquad \langle\Phi_+,\Phi_-\rangle_\mathcal{H}=0.
\]
\end{proposition}

\begin{proof}
By Lemma~\ref{lem:pre-self-inner},
\[
  \|\Phi_{+,X}-\Phi_{+,X'}\|_\mathcal{H}^2\;=\;\langle\Phi_{+,X},\Phi_{+,X}\rangle+\langle\Phi_{+,X'},\Phi_{+,X'}\rangle-2\Re\langle\Phi_{+,X},\Phi_{+,X'}\rangle\;\to\;1+1-2\cdot 1\;=\;0
\]
as $\min(X,X')\to\infty$, with the explicit rate $4/\min(X-X_0,X'-X_0)$. By completeness of $\mathcal{H}$ the limit $\Phi_+$ exists. The limit of $\langle\Phi_{+,X},\Phi_{+,X'}\rangle$ as both $X,X'\to\infty$ is $1$, so $\|\Phi_+\|_\mathcal{H}^2=1$. The analogous argument gives $\Phi_-$ and $\|\Phi_-\|_\mathcal{H}^2=1$. For orthogonality, by Lemma~\ref{lem:pre-self-inner} $\langle\Phi_{+,X},\Phi_{-,X'}\rangle_\mathcal{H}\to 0$ as $\min(X,X')\to\infty$, and continuity of the inner product on $\mathcal{H}$ transfers this to $\langle\Phi_+,\Phi_-\rangle_\mathcal{H}=0$.
\end{proof}

\begin{definition}[The complex orthogonal measure]\label{def:limit-measure}
For every Borel set $A\subset\mathbb{R}$ set
\[
  \Phi(A)\;:=\;\mathbf{1}_A(+1)\,\Phi_+\;+\;\mathbf{1}_A(-1)\,\Phi_-\;\in\;\mathcal{H}.
\]
Equivalently, $\Phi(A)=\lim_{X\to\infty}\Phi_X(A)$ in $\mathcal{H}$-norm, with the limit existing by Proposition~\ref{prop:cauchy}.
\end{definition}

\section{Orthogonal-scattering and intensity}\label{sec:scatter}

\begin{proposition}[Orthogonal scattering]\label{prop:orthogonal}
For every pair of Borel sets $A,B\subset\mathbb{R}$,
\[
  \langle\Phi(A),\Phi(B)\rangle_\mathcal{H}\;=\;\nu(A\cap B),\qquad \nu=\delta_{-1}+\delta_{+1}.
\]
In particular, if $A\cap B=\emptyset$ then $\langle\Phi(A),\Phi(B)\rangle_\mathcal{H}=0$, so $\Phi$ is orthogonally scattered, with intensity measure $\nu$.
\end{proposition}

\begin{proof}
Expand using Definition~\ref{def:limit-measure}:
\[
  \langle\Phi(A),\Phi(B)\rangle_\mathcal{H}\;=\;\mathbf{1}_A(+1)\mathbf{1}_B(+1)\langle\Phi_+,\Phi_+\rangle+\mathbf{1}_A(-1)\mathbf{1}_B(-1)\langle\Phi_-,\Phi_-\rangle+\text{(cross terms)}.
\]
Proposition~\ref{prop:cauchy} gives $\langle\Phi_\pm,\Phi_\pm\rangle=1$ and $\langle\Phi_+,\Phi_-\rangle=0$, so the expression collapses to $\mathbf{1}_A(+1)\mathbf{1}_B(+1)+\mathbf{1}_A(-1)\mathbf{1}_B(-1)=\nu(A\cap B)$.
\end{proof}

\begin{corollary}[$\sigma$-additivity and finiteness]\label{cor:sigma}
The set function $A\mapsto\Phi(A)\in\mathcal{H}$ is $\sigma$-additive in the $\mathcal{H}$-norm topology and satisfies $\sum_n\|\Phi(A_n)\|_\mathcal{H}^2=\nu(\bigsqcup A_n)$ for every countable disjoint family $(A_n)$.
\end{corollary}

\begin{proof}
If $(A_n)_{n\ge 1}$ is a countable disjoint family of Borel sets, only those $A_n$ containing at least one of $\pm 1$ can contribute. Up to relabelling there are at most two such sets, $A_{n_-}$ containing $-1$ and $A_{n_+}$ containing $+1$, and
\[
  \Phi\Big(\bigsqcup_n A_n\Big)\;=\;\Phi_-\mathbf{1}_{\{\exists n:-1\in A_n\}}+\Phi_+\mathbf{1}_{\{\exists m:+1\in A_m\}}\;=\;\sum_n\Phi(A_n),
\]
the series having at most two nonzero terms. The Parseval identity follows from orthonormality of $\{\Phi_-,\Phi_+\}$.
\end{proof}

\section{Spectral representation of \texorpdfstring{$H$}{H}: Fubini verification}\label{sec:fubini}

For $y\in\mathbb{R}$ define the $\mathcal{H}$-valued integral against a bounded Borel function $f:\mathbb{R}\to\mathbb{C}$ by
\[
  \int_{\mathbb{R}} f(\xi)\,d\Phi(\xi)\;:=\;f(+1)\,\Phi_+\;+\;f(-1)\,\Phi_-\;\in\;\mathcal{H}.
\]
This is the standard integral against a $\mathcal{H}$-valued orthogonally-scattered measure supported on the two-point set $\{-1,+1\}$, evaluated as a Riemann sum with exactly two terms.

\begin{theorem}[Spectral representation]\label{thm:spectral}
For every $y\in\mathbb{R}$,
\[
  \widetilde{H_y}\;=\;\int_{\mathbb{R}} e^{i\xi y}\,d\Phi(\xi)\;=\;e^{-iy}\Phi_-+e^{+iy}\Phi_+\qquad\text{as an identity in $\mathcal{H}$.}
\]
\end{theorem}

\begin{proof}
We prove the identity by showing that the difference has zero inner product with every generator $\widetilde{H_s}$, $s\in\mathbb{R}$, and then with itself, so that its $\mathcal{H}$-norm is zero.

\emph{Step 1: Explicit computation of $\langle e^{-iy}\Phi_-+e^{+iy}\Phi_+,\widetilde{H_s}\rangle_\mathcal{H}$ via Fubini.}

Using the Bochner integrals defining $\Phi_{\pm,X}$,
\begin{align*}
  \langle \Phi_{+,X},\widetilde{H_s}\rangle_\mathcal{H}
  &\;=\;\Big\langle\frac{1}{X-X_0}\int_{X_0}^X e^{-ix}\widetilde{H_x}\,dx,\,\widetilde{H_s}\Big\rangle_\mathcal{H}\\
  &\;=\;\frac{1}{X-X_0}\int_{X_0}^X e^{-ix}\,\langle\widetilde{H_x},\widetilde{H_s}\rangle_\mathcal{H}\,dx\qquad\text{(Fubini / Bochner linearity)}\\
  &\;=\;\frac{1}{X-X_0}\int_{X_0}^X e^{-ix}\bigl[e^{-i(x-s)}+e^{+i(x-s)}\bigr]dx\\
  &\;=\;\frac{e^{+is}}{X-X_0}\int_{X_0}^X e^{-2ix}dx\;+\;\frac{e^{-is}}{X-X_0}\int_{X_0}^X 1\,dx\\
  &\;=\;\frac{e^{+is}}{-2i(X-X_0)}\bigl(e^{-2iX}-e^{-2iX_0}\bigr)\;+\;e^{-is}.
\end{align*}
Passing to the limit $X\to\infty$ (the first term vanishes as $1/(X-X_0)$, and Proposition~\ref{prop:cauchy} transports $\Phi_{+,X}\to\Phi_+$ under the continuous linear functional $\langle\cdot,\widetilde{H_s}\rangle_\mathcal{H}$) gives
\[
  \langle\Phi_+,\widetilde{H_s}\rangle_\mathcal{H}\;=\;e^{-is}.
\]
The analogous computation for $\Phi_{-,X}$ (swap $e^{-ix}$ for $e^{+ix}$ in the kernel; the non-vanishing term is the one with matching phase) yields $\langle\Phi_-,\widetilde{H_s}\rangle_\mathcal{H}=e^{+is}$.

Therefore
\[
  \langle e^{-iy}\Phi_-+e^{+iy}\Phi_+,\widetilde{H_s}\rangle_\mathcal{H}\;=\;e^{-iy}\cdot e^{+is}+e^{+iy}\cdot e^{-is}\;=\;e^{-i(y-s)}+e^{+i(y-s)}\;=\;C(y-s).
\]

\emph{Step 2: Matching inner products.} By the definition of the time-average inner product,
\[
  \langle\widetilde{H_y},\widetilde{H_s}\rangle_\mathcal{H}\;=\;C(y-s).
\]
So for every $s\in\mathbb{R}$,
\[
  \langle\widetilde{H_y}-e^{-iy}\Phi_--e^{+iy}\Phi_+,\widetilde{H_s}\rangle_\mathcal{H}\;=\;0.
\]

\emph{Step 3: Self-inner-product computation.} By orthonormality of $\{\Phi_-,\Phi_+\}$ in $\mathcal{H}$,
\[
  \|e^{-iy}\Phi_-+e^{+iy}\Phi_+\|_\mathcal{H}^2\;=\;|e^{-iy}|^2\|\Phi_-\|^2+|e^{+iy}|^2\|\Phi_+\|^2\;=\;2\;=\;C(0)\;=\;\|\widetilde{H_y}\|_\mathcal{H}^2.
\]
Step~1 applied to $s=y$ gives
\[
  \langle e^{-iy}\Phi_-+e^{+iy}\Phi_+,\widetilde{H_y}\rangle_\mathcal{H}\;=\;e^{-iy}\cdot e^{+iy}+e^{+iy}\cdot e^{-iy}\;=\;2,
\]
and the conjugate-symmetric inner product satisfies $\langle\widetilde{H_y},e^{-iy}\Phi_-+e^{+iy}\Phi_+\rangle_\mathcal{H}=\overline{2}=2$. Therefore
\[
  \|\widetilde{H_y}-e^{-iy}\Phi_--e^{+iy}\Phi_+\|_\mathcal{H}^2\;=\;\|\widetilde{H_y}\|^2-2\Re(2)+\|e^{-iy}\Phi_-+e^{+iy}\Phi_+\|^2\;=\;2-4+2\;=\;0.
\]
Hence $\widetilde{H_y}=e^{-iy}\Phi_-+e^{+iy}\Phi_+$ in $\mathcal{H}$, which is exactly $\int_{\mathbb{R}}e^{i\xi y}\,d\Phi(\xi)$ by the definition of the integral against $\Phi$.
\end{proof}

\begin{corollary}[Spectral representation of the Hardy $Z$-function]\label{cor:Z}
For every $t\ge T_0$,
\[
  \widetilde{Z_t}\;:=\;\sqrt{\theta'(t)}\,\widetilde{H_{\theta(t)}}\;=\;\sqrt{\theta'(t)}\int_{\mathbb{R}} e^{i\xi\theta(t)}\,d\Phi(\xi)\;=\;\sqrt{\theta'(t)}\bigl(e^{-i\theta(t)}\Phi_-+e^{+i\theta(t)}\Phi_+\bigr)
\]
as an identity in $\mathcal{H}$.
\end{corollary}

\begin{proof}
$H(x)=Z(\theta^{-1}(x))/\sqrt{\theta'(\theta^{-1}(x))}$, so $Z(t)=\sqrt{\theta'(t)}H(\theta(t))$, and Theorem~\ref{thm:spectral} at $y=\theta(t)$ gives the stated identity on scaling by $\sqrt{\theta'(t)}$.
\end{proof}

\section{Isometry and reality}\label{sec:isometry}

\begin{proposition}[Isometry onto $L^2(\nu)$]\label{prop:isometry}
The map
\[
  V:L^2(\mathbb{R},\nu)\longrightarrow\mathcal{H},\qquad V(f)\;:=\;\int_{\mathbb{R}} f(\xi)\,d\Phi(\xi)\;=\;f(+1)\Phi_++f(-1)\Phi_-,
\]
is a surjective linear isometry onto $\overline{\mathrm{span}}\{\Phi_+,\Phi_-\}\subseteq\mathcal{H}$. In particular $\dim\mathcal{H}=2$.
\end{proposition}

\begin{proof}
$L^2(\mathbb{R},\nu)$ is $2$-dimensional with orthonormal basis $\{\mathbf{1}_{\{+1\}},\mathbf{1}_{\{-1\}}\}$. $V$ sends this basis to $\{\Phi_+,\Phi_-\}$, which is orthonormal in $\mathcal{H}$ by Proposition~\ref{prop:cauchy}. So $V$ is an isometry onto its image. The image contains every $\widetilde{H_s}=e^{-is}\Phi_-+e^{+is}\Phi_+=V(e^{-is}\mathbf{1}_{\{-1\}}+e^{+is}\mathbf{1}_{\{+1\}})$ by Theorem~\ref{thm:spectral}, so it contains the dense subspace $\mathcal{E}/\mathcal{N}$, hence equals $\mathcal{H}$.
\end{proof}

\begin{corollary}[Reality factorization]\label{cor:reality}
Under the isometry $V$ of Proposition~\ref{prop:isometry}, the generator $\widetilde{H_y}$ corresponds to the $L^2(\nu)$-function $\xi\mapsto e^{i\xi y}$, i.e.\ $V^{-1}(\widetilde{H_y})(\xi)=e^{i\xi y}$ for $\xi\in\{-1,+1\}$. The two values of this function on the support of $\nu$ are $e^{+iy}$ and $e^{-iy}$, which are complex conjugates of one another. This conjugate symmetry on the $\nu$-support is the precise expression in the spectral representation of the fact that $H$ is real-valued.
\end{corollary}

\begin{proof}
Immediate from Theorem~\ref{thm:spectral}: $V(\xi\mapsto e^{i\xi y})=e^{-iy}\Phi_-+e^{+iy}\Phi_+=\widetilde{H_y}$, and $\overline{e^{+iy}}=e^{-iy}$.
\end{proof}

\section{Summary}

The complex orthogonal measure $\Phi$ of the warped Hardy signal $H$ is constructed explicitly by the Cram\'er-type pre-limit formula of Definition~\ref{def:prelimit-band}: 
\[
  \Phi_{\pm,X}\;=\;\frac{1}{X-X_0}\int_{X_0}^X e^{\mp ix}\widetilde{H_x}\,dx.
\]
These are Cauchy in the time-average Hilbert space $\mathcal{H}$ by the direct inner-product computation of Lemma~\ref{lem:pre-self-inner} (which reduces to evaluating two one-variable oscillatory integrals, each bounded by $1$), and their limits $\Phi_\pm$ are orthonormal in $\mathcal{H}$. The limit measure $\Phi(A)=\mathbf{1}_A(+1)\Phi_++\mathbf{1}_A(-1)\Phi_-$ is $\sigma$-additive, orthogonally scattered, and has intensity $\nu=\delta_{-1}+\delta_{+1}$. Substituting $\Phi_{\pm,X}$ into $\langle\cdot,\widetilde{H_s}\rangle_\mathcal{H}$ and applying Fubini/Bochner linearity to interchange the inner product with the Bochner integral gives, after evaluating the two resulting one-variable oscillatory integrals and passing to the limit,
\[
  \langle\Phi_\pm,\widetilde{H_s}\rangle_\mathcal{H}\;=\;e^{\mp is},
\]
and matching inner products against the dense family $\{\widetilde{H_s}\}$ (together with the matching self-norm) closes the spectral representation
\[
  \widetilde{H_y}\;=\;\int_{\mathbb{R}} e^{i\xi y}\,d\Phi(\xi)\;=\;e^{-iy}\Phi_-+e^{+iy}\Phi_+.
\]
The Hardy $Z$-function inherits the representation via the scalar lift $Z(t)=\sqrt{\theta'(t)}H(\theta(t))$.

\end{document}
